\section{The remaining determinant estimate and the existence argument}
\label{dfn-next}

The leading correction is now an actual functional,
\[
 D_{M,R}=\Delta_{M,R}+\mathcal Q_{\mathrm{det},M,R},
\]
with an exact Hessian trace formula and an exact finite
composition law. In a jointly compatible infinite system its
Maxwell coefficient satisfies
\[
 \beta_{MRL}(D_{M,RL})
   =\beta_{MR}(D_{M,R})+\beta_{MRL}(D_{MR,L}).
\]
These identities supply the algebra for accumulated drift;
they do not determine its value or sign.

The exact gauge and coarea bridge makes the comparison
operators concrete. With its stated finite coordinates,
\[
 D_{M,K}=\mathcal N_0\!\left[
    \log\det G_b-\frac12\log\det Q_b
                 -\frac12\log\det(CQ_b^{-1}C^T)\right],
\]
where \(\mathcal N_0\) subtracts the value at the identity.
The inverse pivot has been accounted for by coarea.
In fact \(CQ_b^{-1}C^T=(q_{\mathrm{eff}}+B)^{-1}\).
Choosing one fixed sufficiently large \(B=aI\), the
coarse term \(\tfrac12\mathcal N_0\log\det(q_{\mathrm{eff}}+aI)\)
has a uniformly bounded Maxwell coefficient by the
proved coarse-action support bounds.
The remaining comparison concerns the fine combination
\[
 \mathcal N_0\!\left[
 \log\det G_b-\frac12\log\det
 \bigl(s''-\lambda\cdot f''+D_bD_b^T+aC^TC\bigr)\right].
\]
In particular, the point-anchor Rayleigh test rules out a
uniform physical mass estimate for the based ghost alone.

A concrete route is to compare the composite determinant from
the bare Wilson action to a background-field vector and ghost
reference at the same cutoff and infrared scale. With every
normalization specified, the desired estimate has the form
\[
 \bigl|\beta(D_{1,R^k})-B_k^{\mathrm{ref}}\bigr|\le C,
 \qquad
 \bigl|B_k^{\mathrm{ref}}-a\,k\bigr|\le C',
 \qquad a>0,
\]
with constants independent of \(k\). Neither displayed bound is
proved in this notebook. The stationary lift, constrained-current
curvature and measure factors must be included in that comparison.
A coefficient quoted for a different effective action does not
establish these estimates for the present one.

The drift hypothesis required for shooting concerns the actual
finite regulators, including arbitrary permitted preceding depths.
An infinite comparison must therefore retain the uniformity and
joint-limit compatibility needed to transfer a positive block
margin to those finite coefficients. The qualitative transfer
proposition supplies that passage once a positive margin holds
in every relevant structural limit; a quantitative
finite-volume convergence rate is unnecessary for this purpose.
No exact finite commutation of Maxwell extraction with classical
transport has been assumed.

Once that drift is supplied, the conditional shooting theorem
closes the finite weak-coupling trajectory issue for this local
map. The subsequent existence argument must still incorporate
all required field histories with their exact measure identities,
prove uniform source estimates, construct a nontrivial continuum
observable hierarchy, and establish its Euclidean symmetry,
positivity and reconstruction properties. These requirements are
those of Notebook 01; the determinant and tuning arguments do not
replace them.

\begin{thebibliography}{9}
\bibitem{priorflow}
\emph{Toward a four-dimensional Yang--Mills construction},
research Notebooks 15--31 in this workspace.
The local minimum, exact constrained integral, common-domain
quantum map, finite induction and classical Maxwell transport
are the principal inputs.
\bibitem{lwbackground}
M.~L\"uscher and P.~Weisz,
\emph{Background field technique and renormalization in lattice
gauge theory}, Nucl. Phys. B \textbf{452} (1995), 213--233.
\href{https://arxiv.org/abs/hep-lat/9504006}{Primary preprint}.
\bibitem{lwmatching}
M.~L\"uscher and P.~Weisz,
\emph{Computation of the relation between the bare lattice
coupling and the MS coupling in SU(N) gauge theories to two
loops}, Nucl. Phys. B \textbf{452} (1995), 234--260.
\href{https://arxiv.org/abs/hep-lat/9505011}{Primary preprint}.
These papers supply background-field formalism and reference
amplitudes to investigate. They are not used here as a proof of
the required comparison with our constrained determinant.
\bibitem{sourceflow}
E.~Sussman and collaborators,
\href{https://github.com/ethanwsussman/Yang--Mills-LLM-experiment}
{Yang--Mills-LLM-experiment}, deposited raw accounts at commit
\texttt{e392ef45013a581664bab76b90867170b30fe10e}.
\end{thebibliography}
\end{document}
